% Fig.: what goes in, what comes out (main model), how it is tested, and the separate after-release model.
\begin{figure*}[t]
\centering
\begin{tikzpicture}[font=\footnotesize, >=Stealth, node distance=4mm,
  src/.style={draw=black!45, fill=black!4, rounded corners=2pt, align=center, text width=2.6cm, minimum height=9mm},
  feat/.style={draw=blue!60!black, fill=blue!6, rounded corners=2pt, align=center, text width=3.7cm, minimum height=9mm},
  stp/.style={draw=black!70, fill=white, rounded corners=2pt, align=center, minimum height=9mm},
  outb/.style={draw=blue!60!black, fill=blue!15, rounded corners=2pt, align=center, minimum height=9mm},
  lab/.style={font=\scriptsize\itshape, text=black!60}]
% ---- inputs
\node[src] (a) {Song audio\\(Spotify Web API)};
\node[src, below=of a] (d) {Release date};
\node[src, below=of d] (h) {Artist name +\\Billboard Hot~100, 1958--};
\node[lab, above=1mm of a] {1. raw inputs};
% ---- features
\node[feat, right=7mm of a] (fa) {13 audio descriptors\\{\scriptsize danceability, energy, loudness, \dots}};
\node[feat, right=7mm of d] (fd) {decade $\rightarrow$ 6 indicators\\{\scriptsize 1960s \dots\ 2010s, one-hot}};
\node[feat, right=7mm of h] (fh) {artist history \mbox{$h(s)=\log(1+n)$}\\{\scriptsize $n$ = its Hot~100 songs in earlier decades}};
\node[lab, above=1mm of fa] {2. features (20 numbers)};
\foreach \x/\y in {a/fa, d/fd, h/fh} \draw[->] (\x) -- (\y);
% ---- model
\node[stp, right=9mm of fd, text width=2.5cm] (rf) {\textbf{Random Forest}\\80 trees, 8 features\\per split, never tuned};
\foreach \x in {fa, fd, fh} \draw[->] (\x.east) -- (rf.west);
\node[lab, above=1mm of rf] {3. model};
\node[outb, right=7mm of rf, text width=1.6cm] (p) {$P(\text{hit})$\\$\in[0,1]$};
\node[outb, right=7mm of p, text width=2.1cm] (dec) {hit if\\$P(\text{hit})\geq 0.5$};
\draw[->] (rf) -- (p);
\draw[->] (p) -- (dec);
\node[lab, above=1mm of p, xshift=6mm] {4. output};
% ---- evaluation band
\node[stp, below=8mm of h.south west, anchor=north west, text width=16.2cm, fill=black!3, draw=black!40] (ev)
 {\textbf{How every number is tested.} 5 folds $\times$ 5 repeats: the model is trained on 80\,\% of the songs and
 scored on the 20\,\% it never saw, 25 times $\rightarrow$ mean accuracy and 95\,\% CI $\rightarrow$ corrected
 $t$-test (Nadeau--Bengio) against the same model without the feature $\rightarrow$ Bonferroni over all
 $\approx$\,60 tests ($p<8\times10^{-4}$) $\rightarrow$ controls: year-matched negatives, genre-matched negatives,
 no artist in both training and test folds.};
% ---- after release
\node[src, below=10mm of ev.south west, anchor=north west, text width=2.6cm] (w) {Chart positions in\\weeks 1, 2, 3 after debut};
\node[feat, right=7mm of w, text width=4.6cm] (wf) {positions, best rank, climb per week,\\acceleration + year, month,\\artist's earlier hits};
\node[stp, right=9mm of wf, text width=2.5cm] (rf2) {Random Forest\\200 trees};
\node[outb, right=7mm of rf2, text width=3.9cm] (p2) {$P(\text{peak in the top 10})$\\\scriptsize separate question: only after release};
\draw[->] (w) -- (wf);
\draw[->] (wf) -- (rf2);
\draw[->] (rf2) -- (p2);
\node[lab, above=1mm of w.north west, anchor=south west] {After release (Sec.~\ref{sec:timeseries}): a second, separate model};
\end{tikzpicture}
\caption{The prediction pipeline. Top: the final model predicts, on release day, whether a song will enter the
Hot~100, from 20 numbers per song. Middle: the protocol under which every gain in this paper is measured. Bottom:
a separate model, used only once a song has charted, predicts whether it will reach the top~10 from its first three
weeks.}
\label{fig:pipeline}
\end{figure*}
